problem:  
**(Bach–Positive Curvature Rigidity Conjecture in Four Dimensions)**  
Let \((M^4,g)\) be a compact, simply connected Riemannian 4‑manifold with **positive sectional curvature**.  
Suppose that \(g\) is **Bach‑flat**, i.e. it satisfies  
\[
B_{ij} = \nabla^k \nabla^\ell W_{ikj\ell} + \tfrac{1}{2} R^{k\ell} W_{ikj\ell} = 0,
\]  
so that \(g\) is a stationary point of the **Weyl energy**
\[
\mathcal{W}(g) = \int_M |W|^2\, d\mathrm{vol}_g
\]
under conformal variations.

**Question:**  
Must such a metric necessarily be **Einstein**, and consequently is it necessarily **isometric to a compact rank–one symmetric space (CROSS)**, i.e. \(S^4\) or \(\mathbb{CP}^2\) with their standard metrics?

Why is it a "good" problem:  
This problem formulates the precise and open question lying at the intersection of *conformal geometry*, *curvature variational theory*, and *positive curvature rigidity*. It is “good” in the following strong senses:

1. **Sharp interface of conformal and curvature analysis:**  
   The Bach–flat condition represents the Euler–Lagrange equation of the conformally invariant \(L^2\)-Weyl functional in dimension four, while positive sectional curvature encodes strict Riemannian positivity. Their interaction is substantially unexplored — it asks whether conformal criticality plus positivity collapses to Einstein rigidity.

2. **Bridges distinct frameworks:**  
   Resolving this would unify two major rigidity theories: (a) *conformal critical metric theory* (Bach–flat geometry) and (b) *sphere theorem–type positivity rigidity*. A proof would connect fourth‑order conformal PDEs with second‑order curvature–operator inequalities.

3. **Analytic and geometric depth:**  
   The Bach equation is an elliptic fourth‑order system invariant under conformal changes. Showing that positive sectional curvature enforces Einstein structure would require a novel combination of *Bochner–Weitzenböck identities*, *curvature operator positivity analyses*, or possibly *Ricci flow continuation from Bach‑flat initial data*—all delicate analytic machinery.

4. **Simplicity and universality:**  
   Easily stated with classical tensors, it isolates an appealing question: *Can positivity of curvature eliminate non‑Einstein critical points of the Weyl energy?* The geometric content is transparent yet conceals formidable analytic challenges.

5. **Truly open and structurally plausible:**  
   No classification of positive‑curvature Bach‑flat metrics is known. Einstein metrics like \(S^4\) and \(\mathbb{CP}^2\) trivially satisfy the hypotheses, but whether they are the only ones remains open. A solution would close a significant gap between conformal geometry and positive curvature rigidity.

Hence, this conjecture is a concise, natural, and genuinely open research‑level question—simple to state, profoundly deep, and integrative across geometric analysis and Riemannian topology.